% ============================================================
% ADAPTIVE TOOLKIT, not a fixed per-chapter template.
% Each chapter uses whichever of these actually serves its content.
% None are mandatory; none appear in a fixed order. The one standing
% expectation (see SKILL.md, Book Design): any chapter that explains a
% mechanism shows it as a picture as well as in words.
% ============================================================

% ------------------------------------------------------------
% Text-level tools
% ------------------------------------------------------------

% Colon headings: "Label: elaboration".
\newcommand{\theading}[2]{%
  \par\Needspace{9\baselineskip}\addvspace{20pt}%
  \noindent{\displayfont\bfseries\large\color{accent}#1:}%
  \hspace{0.4em}{\displayfont\large\color{ink}#2}%
  \par\addvspace{8pt}\nopagebreak
}
\newcommand{\tlabel}[1]{{\displayfont\bfseries\large\color{accent}#1}\par\addvspace{6pt}\nopagebreak}

% Sidenote: a short aside (a definition, a date, a small joke) set at the
% foot of the page, numbered in the accent color.
\renewcommand{\thefootnote}{\arabic{footnote}}
\makeatletter
\renewcommand{\@makefnmark}{\hbox{\textsuperscript{\color{accent}\bfseries\@thefnmark}}}
\renewcommand{\footnoterule}{\vspace*{-3pt}{\color{accent!50}\hrule width 1.2in height 0.5pt}\vspace*{6pt}}
\makeatother
\renewcommand{\footnotelayout}{\itshape\color{ink!85}}
\setlength{\footnotesep}{10pt}
\newcommand{\sidenote}[1]{\footnote{#1}}

% Drop cap for a chapter's first paragraph: \opening{W}{hen we type}
\renewcommand{\LettrineFontHook}{\displayfont\bfseries\color{accent}}
\setlength{\DefaultFindent}{2pt}
\setlength{\DefaultNindent}{0pt}
\newcommand{\opening}[2]{\lettrine[lines=3,loversize=0.05]{#1}{#2}}

% Epigraph: an optional opening frame.
\newcommand{\epigraph}[2]{%
  \begin{flushright}\displayfont\itshape\color{ink!80}
    \parbox{0.8\textwidth}{\raggedleft #1\par\smallskip
    \normalfont\upshape\small\color{accent}#2}
  \end{flushright}\bigskip
}

% Depth tag: one per chapter, right under the chapter title:
% \depthtag{core} | {practical} | {deep dive} | {specialized}
\newcommand{\depthtag}[1]{%
  \par\noindent\hfill\tikz[baseline=(t.base)]\node[
    rounded rectangle, fill=accent!12, text=accent!80!black,
    font=\scriptsize\bfseries\addfontfeatures{LetterSpace=10}, inner xsep=8pt, inner ysep=3pt](t){\MakeUppercase{#1}};%
  \par\vspace{6pt}}

% Full-width block: kept for wide tables and figures; with symmetric
% margins it simply centers its content on the text block.
\newenvironment{fullwidth}{\par\centering}{\par}

% ------------------------------------------------------------
% Boxes. Each has one job; the title tells the reader which.
% ------------------------------------------------------------

% Shared skeleton: a soft panel with a colored spine and a small label.
\tcbset{
  fz/.style={
    breakable, enhanced, boxrule=0pt, frame hidden, arc=2pt,
    left=12pt, right=10pt, top=8pt, bottom=8pt,
    before skip=12pt, after skip=12pt,
    fonttitle=\displayfont\bfseries\small,
    attach title to upper, after title={\par\smallskip},
  },
}

% The mental model: the one picture to keep when the details fade.
\newtcolorbox{mentalmodel}[1][]{fz, colback=accent!6,
  borderline west={3pt}{0pt}{accent},
  coltitle=accent, title={The mental model\ifx\relax#1\relax\else: #1\fi}}

% Course refresher: the short recap the audience rule promises.
% \begin{refresher}{CSE321}{Virtual memory} ... \end{refresher}
\newtcolorbox{refresher}[2]{fz, colback=cool!6,
  borderline west={3pt}{0pt}{cool}, coltitle=cool,
  title={Quick refresher: #2 \hfill{\normalfont\scriptsize from #1}}}

% Back in time: the history behind a design, with its year.
\newtcolorbox{backintime}[1]{fz, colback=gold!10,
  borderline west={3pt}{0pt}{gold!80!black}, coltitle=gold!45!black,
  title={Back in time: #1}}

% Predict first: a question the reader answers before reading on.
\newtcolorbox{predict}{fz, colback=white,
  borderline={0.6pt}{0pt}{accent!60,dashed}, coltitle=accent,
  title={Predict first}}

% Myth vs. reality: a common belief, then what actually happens.
% \begin{myth}{Containers are lightweight virtual machines.} ... \end{myth}
\newtcolorbox{myth}[1]{fz, colback=paper,
  borderline west={3pt}{0pt}{ink!60}, coltitle=ink,
  title={Myth: \normalfont\itshape #1}, after title={\par\smallskip\textbf{\color{accent}Reality.}\space}}

% Break it on purpose: the deliberate-failure moment.
\newtcolorbox{breakage}{fz, colback=warn!5,
  borderline west={3pt}{0pt}{warn}, coltitle=warn,
  title={Break it on purpose}}

% Trace it: the part-end exercise, one operation through every layer.
\newtcolorbox{traceit}[1]{fz, unbreakable, colback=cool!5,
  borderline={0.8pt}{0pt}{cool}, coltitle=cool, left=14pt,
  title={Trace it: \normalfont\itshape #1}}

% FAQ: one per real question, wherever it is relevant.
\newtcolorbox{faqentry}[1]{fz, colback=accent!4,
  borderline west={1.5pt}{0pt}{accent!60}, coltitle=accent!60!black,
  title={#1}}

% ------------------------------------------------------------
% Code, terminals, algorithms
% ------------------------------------------------------------

% Terminal: real command, real output.
\lstdefinestyle{termstyle}{
  basicstyle=\ttfamily\small\color{codeink},
  backgroundcolor=\color{codebg},
  frame=l, framerule=2.5pt, rulecolor=\color{accent},
  framesep=10pt, xleftmargin=12pt, breaklines=true,
  columns=fullflexible, keepspaces=true, showstringspaces=false,
  aboveskip=12pt, belowskip=12pt
}
\lstnewenvironment{terminal}{\lstset{style=termstyle}}{}

% Source code: \begin{code}{Python} ... \end{code}, with gentle syntax color.
\lstdefinestyle{codestyle}{
  basicstyle=\ttfamily\small\color{codeink},
  backgroundcolor=\color{paper},
  keywordstyle=\color{accent}\bfseries,
  commentstyle=\color{ink!55}\itshape,
  stringstyle=\color{cool},
  numbers=left, numberstyle=\tiny\color{ink!40}, numbersep=8pt,
  frame=none, xleftmargin=16pt, framesep=6pt,
  breaklines=true, columns=fullflexible, keepspaces=true,
  showstringspaces=false, aboveskip=12pt, belowskip=12pt
}
\lstdefinelanguage{JavaScript}{
  keywords={const,let,var,function,return,if,else,for,while,import,export,
            from,async,await,new,class,true,false,null,undefined,require},
  sensitive=true, comment=[l]{//}, morecomment=[s]{/*}{*/},
  morestring=[b]', morestring=[b]", morestring=[b]`
}
\lstnewenvironment{code}[1]{\lstset{style=codestyle,language=#1}}{}

% Algorithm: numbered pseudocode in a quiet frame.
% \begin{algo}{How the shell resolves a command name} ... \end{algo}
\lstdefinelanguage{pseudo}{
  keywords={for,each,in,if,then,else,return,while,do,end,and,or,not,
            procedure,function,fail,break},
  sensitive=false, comment=[l]{//}
}
\newtcblisting{algo}[1]{fz, listing only, colback=white,
  borderline north={0.8pt}{0pt}{ink}, borderline south={0.8pt}{0pt}{ink},
  coltitle=ink, left=4pt, title={Algorithm: \normalfont #1},
  listing options={style=codestyle, language=pseudo, backgroundcolor=\color{white},
    keywordstyle=\bfseries\color{ink}, aboveskip=0pt, belowskip=0pt}}

% Comparison tables: booktabs, no vertical rules, generous rows.
\newcommand{\comparerow}{\\[2pt]}
\renewcommand{\arraystretch}{1.3}

% ------------------------------------------------------------
% Diagram vocabulary (TikZ). Every figure in the book draws from these
% styles so the whole book reads as one visual system.
% ------------------------------------------------------------
\tikzset{
  fz/.style={font=\small, >={Stealth[length=6pt]}, line width=0.7pt},
  box/.style={draw=ink!70, fill=paper, rounded corners=3pt,
              minimum height=8mm, inner xsep=8pt, align=center},
  proc/.style={box, draw=accent, fill=accent!8},          % a running process
  kern/.style={box, draw=cool, fill=cool!10},             % kernel / system side
  file/.style={box, draw=ink!60, fill=white, sharp corners,
               tape, tape bend top=none},                 % a file on disk
  data/.style={box, draw=gold!80!black, fill=gold!12},    % data, a value, a variable
  note/.style={font=\footnotesize\itshape, text=ink!75, align=left},
  flow/.style={->, draw=ink!75},
  call/.style={->, draw=accent, line width=0.9pt},
  reply/.style={->, draw=cool, dashed},
  boundary/.style={draw=cool, dashed, line width=0.8pt},
  layer/.style={rounded corners=2pt, minimum height=7mm, inner xsep=6pt,
                font=\small\bfseries, text=white},
  stepdot/.style={circle, fill=accent, text=white, font=\scriptsize\bfseries,
               inner sep=1.5pt, minimum size=4.5mm},
}

% \layerstack{width}{{Label one},{Label two},...}: a stack of layers,
% first item on top, colored from cool (top) to accent (bottom).
\newcommand{\layerstack}[2]{%
  \begin{tikzpicture}[fz]
    \foreach \lbl [count=\i] in {#2} {\global\let\fzlayercount\i}
    \foreach \lbl [count=\i] in {#2} {
      \pgfmathsetmacro{\mix}{100*(\i-1)/max(1,\fzlayercount-1)}
      \node[layer, fill=accent!\mix!cool, minimum width=#1]
        at (0,-\i*8.5mm) {\lbl};
    }
  \end{tikzpicture}%
}

% ------------------------------------------------------------
% Chapter, part, section, and page furniture
% ------------------------------------------------------------

\newcommand{\fzchaptitle}[1]{%
  \begin{minipage}{\fullw}\raggedright\displayfont\bfseries
    \fontsize{25}{30}\selectfont\color{ink}#1\par\vspace{10pt}
    {\color{accent}\rule{\fullw}{1.2pt}}
  \end{minipage}}

% Chapter opening: small spaced "Chapter" label and a big numeral, the
% title across the full width, an accent rule under it.
\titleformat{\chapter}[display]
  {\normalfont}
  {\makebox[\fullw][l]{{\color{accent}\addfontfeatures{LetterSpace=18}\small\bfseries CHAPTER}\hfill
     {\displayfont\bfseries\fontsize{60}{60}\selectfont\color{accent!28}\thechapter}}}
  {6pt}
  {\fzchaptitle{#1}}
\titleformat{name=\chapter,numberless}[display]
  {\normalfont}{}{0pt}
  {\fzchaptitle{#1}}
\titlespacing*{\chapter}{0pt}{0pt}{30pt}

% Part opening: a full-bleed page. Big numeral, the part title, and a
% map of the whole book with this part lit up, so every part opens by
% showing where we are on the journey. \bookparts is written by
% generate_stubs.py into _generated_structure.tex.
\providecommand{\bookparts}{}
\newcommand{\partmap}{%
  \begin{tikzpicture}[font=\scriptsize]
    \foreach \t [count=\i] in \bookparts {
      \ifnum\i=\value{part}
        \node[anchor=west, text=white, font=\scriptsize\bfseries]
          at (0,-\i*4.3mm) {\makebox[2.4em][l]{\uppercase\expandafter{\romannumeral\i}}\t};
        \fill[gold] (-0.35,-\i*4.3mm) circle (1.3pt);
      \else
        \node[anchor=west, text=white!65!accent] at (0,-\i*4.3mm) {\makebox[2.4em][l]{\uppercase\expandafter{\romannumeral\i}}\t};
      \fi
    }
  \end{tikzpicture}%
}
\titleformat{\part}[block]
  {\normalfont}{}{0pt}
  {\thispagestyle{empty}%
   \begin{tikzpicture}[remember picture, overlay]
     \fill[accent] (current page.north west) rectangle (current page.south east);
     \fill[gold] ([xshift=0.8in,yshift=-3.85in]current page.north west) rectangle
       ([xshift=1.6in,yshift=-3.9in]current page.north west);
     \node[anchor=north west, text=white!80!accent, font=\displayfont\bfseries\fontsize{110}{100}\selectfont]
       at ([xshift=0.75in,yshift=-0.7in]current page.north west) {\Roman{part}};
     \node[anchor=north west, text width=5.6in, align=left, text=white,
           font=\displayfont\bfseries\fontsize{30}{35}\selectfont]
       at ([xshift=0.78in,yshift=-2.55in]current page.north west) {#1};
     \node[anchor=north west, text=white!80!accent, font=\small\addfontfeatures{LetterSpace=18}\bfseries]
       at ([xshift=0.8in,yshift=-2.25in]current page.north west) {PART \Roman{part}};
     \node[anchor=north west] at ([xshift=0.8in,yshift=-4.4in]current page.north west) {\partmap};
   \end{tikzpicture}}
\titlespacing*{\part}{0pt}{0pt}{0pt}

\titleformat{\section}{\displayfont\bfseries\Large\color{ink}}{\color{accent}\thesection}{8pt}{#1}
\titleformat{\subsection}{\displayfont\bfseries\large\color{ink}}{\color{accent}\thesubsection}{6pt}{#1}

% Running heads span the full width, so the rule ties text and margin together.
\pagestyle{fancy}
\fancyhf{}
\fancyhead[L]{\footnotesize\color{ink!65}\addfontfeatures{LetterSpace=8}\MakeUppercase{\leftmark}}
\fancyhead[R]{\footnotesize\bfseries\color{accent}\thepage}
\renewcommand{\headrulewidth}{0.4pt}
\renewcommand{\headrule}{\hbox to\headwidth{\color{accent!40}\leaders\hrule height \headrulewidth\hfill}}
\renewcommand{\chaptermark}[1]{\markboth{\ifnum\value{chapter}>0 \thechapter\quad\fi #1}{}}
\fancypagestyle{plain}{\fancyhf{}%
  \renewcommand{\headrulewidth}{0pt}%
  \fancyfoot[R]{\footnotesize\bfseries\color{accent}\thepage}}

% Contents: parts lead in the accent color, chapters are quiet ink with
% dotted leaders, so the page reads as a map rather than a wall of bold.
\usepackage{titletoc}
\titlecontents{part}[0pt]
  {\addvspace{14pt}\displayfont\bfseries\color{accent}}
  {\makebox[2.2em][l]{\thecontentslabel}}{}
  {\hfill\contentspage}[\addvspace{4pt}]
\titlecontents{chapter}[2.2em]
  {\addvspace{1.5pt}\small\color{ink}}
  {\contentslabel[\color{accent!80}\thecontentslabel]{2.2em}}{}
  {\titlerule*[0.7pc]{.}\contentspage}
